%=============================================================================
%  Chapitre 8 — Sprint 5 : Compagnon IA Melo, Scénarios et Surveillance Énergétique
%=============================================================================
\chapter{Sprint 5 --- Compagnon IA (Melo), Automatisation par Scénarios et Surveillance Énergétique}

\begin{tcolorbox}[sprintbox,title={Sprint 5 --- Compagnon IA \& Automatisation}]
\textbf{Durée :} 2 semaines \quad
\textbf{Points de Story :} 34 \quad
\textbf{Objectif :} Livrer le compagnon IA Melo pour le contrôle des appareils en langage
naturel, l'automatisation temporelle via node-cron, et un tableau de bord de surveillance
de la consommation énergétique.
\end{tcolorbox}

\section{Backlog Sprint 5}

\begin{center}
\begin{longtable}{p{0.07\textwidth} p{0.54\textwidth} p{0.08\textwidth} p{0.12\textwidth}}
\caption{Décomposition des User Stories Sprint 5}\label{tab:sprint5_backlog}\\
\toprule
\textbf{ID} & \textbf{User Story} & \textbf{PS} & \textbf{Statut} \\
\midrule
\endfirsthead
\toprule
\textbf{ID} & \textbf{User Story} & \textbf{PS} & \textbf{Statut} \\
\midrule
\endhead
\bottomrule
\endfoot
US-23 & Compagnon IA (Melo) pour contrôle d'appareils en langage naturel & 13 & \done Terminé \\
US-24 & Scénarios d'automatisation temporels avec node-cron & 8 & \done Terminé \\
US-25 & Graphiques de consommation énergétique horaire & 8 & \done Terminé \\
US-26 & Historique et statut des demandes de permission & 5 & \done Terminé \\
\end{longtable}
\end{center}

\section{Compagnon IA Melo}

\subsection{Architecture}

Melo est construit sur l'\textbf{API Groq}, un service d'inférence cloud offrant des temps
de réponse LLM inférieurs à la seconde. Le backend expose un seul endpoint
\api{POST /api/companion/chat} qui :

\begin{enumerate}[leftmargin=1.5cm,itemsep=4pt]
  \item Enrichit le message utilisateur avec un prompt système structuré contenant la liste
        des appareils courants, les lectures des capteurs et les permissions utilisateur.
  \item Envoie l'historique de conversation à l'API Groq avec le modèle
        \texttt{llama-3.3-70b-versatile}.
  \item Analyse la réponse LLM pour détecter les \textbf{blocs d'action} (commandes JSON structurées).
  \item Exécute les actions validées via le même pipeline de commande d'appareils que l'API REST.
  \item Retourne la réponse LLM et les résultats des actions au frontend.
\end{enumerate}

\subsection{Construction du Prompt Système}

\begin{lstlisting}[style=codeblock,language=JavaScript,caption={Construction du Prompt Syst\`{e}me Melo (backend/routes/companion.js)}]
const buildSystemPrompt = (user, devices, houseState) => {
  const accessibleDevices = user.role === 'admin'
    ? devices
    : devices.filter(d =>
        d.room === user.room ||
        user.permissions.includes(`room:${d.room.toLowerCase().replace(' ', '-')}`) ||
        user.permissions.includes('global:controls')
      );

  return `
Tu es Melo, un assistant maison intelligente intelligent et amical.
Tu aides les utilisateurs a controler leurs appareils en langage naturel.

UTILISATEUR COURANT :
- Nom : ${user.name}
- Role : ${user.role}
- Piece assignee : ${user.room || 'Toutes les pieces (admin)'}

APPAREILS ACCESSIBLES (${accessibleDevices.length} au total) :
${accessibleDevices.map(d =>
  `- ${d.name} (${d.type}) dans ${d.room} | topic: ${d.mqttTopic} | etat: ${d.state}`
).join('\n')}

ENVIRONNEMENT :
- Niveau gaz : ${houseState?.gas?.level?.toFixed(0) ?? 'N/A'} PPM
- Temperature : ${houseState?.temperature ?? 'N/A'}C
- Humidite : ${houseState?.humidity ?? 'N/A'}%
- Scene courante : ${houseState?.lightingMode ?? 'Aucune'}

CAPACITES :
- Pour controler un appareil, reponds avec un bloc ACTION :
  [ACTION: {"deviceId":"<id>","command":{"state":"ON"|"OFF","brightness":0-100}}]
- Pour activer une scene : [SCENE: "Morning"|"Dinner"|"Cinematic"|"Sleep"]
- Pour des requetes d'information, reponds de maniere conversationnelle.
- Ne controle JAMAIS des appareils hors de la liste accessible.
- Si tu ne peux pas faire quelque chose (permission refusee), explique poliment.
`.trim();
};
\end{lstlisting}

\subsection{Implémentation de l'Endpoint Chat}

\begin{lstlisting}[style=codeblock,language=JavaScript,caption={Endpoint Chat Compagnon (backend/routes/companion.js)}]
router.post('/chat', protect, async (req, res) => {
  const { messages } = req.body;   // Historique de conversation du frontend

  const devices    = await Device.find({ houseCode: req.user.houseCode });
  const houseState = await HouseState.findOne({ houseCode: req.user.houseCode });
  const sysPrompt  = buildSystemPrompt(req.user, devices, houseState);

  // Appel API Groq
  const completion = await groq.chat.completions.create({
    model: 'llama-3.3-70b-versatile',
    messages: [
      { role: 'system', content: sysPrompt },
      ...messages
    ],
    temperature: 0.3,              // Basse temperature pour des commandes deterministes
    max_tokens: 1024
  });

  const reply = completion.choices[0].message.content;

  // Analyser et executer les blocs ACTION
  const actionMatches = [...reply.matchAll(/\[ACTION: ({.*?})\]/g)];
  const results = [];
  for (const match of actionMatches) {
    try {
      const { deviceId, command } = JSON.parse(match[1]);
      const device = devices.find(d => d._id.toString() === deviceId);
      if (device) {
        await executeDeviceCommand(device, command, mqttClient, io, req.user);
        results.push({ deviceId, success: true });
      }
    } catch { /* Ignorer les blocs d'action malformes */ }
  }

  // Analyser les blocs SCENE
  const sceneMatch = reply.match(/\[SCENE: "([^"]+)"\]/);
  if (sceneMatch) {
    await activateScene(sceneMatch[1], req.user.houseCode, mqttClient, io);
  }

  res.json({ reply, actionsExecuted: results.length });
});
\end{lstlisting}

\subsection{Composant Frontend Melo}

L'interface Melo est un widget de chat flottant rendu dans le coin inférieur droit de
l'application. Elle inclut une mascotte 3D animée (\texttt{Melo3D}) construite avec
React Three Fiber et animée avec Framer Motion.

\begin{lstlisting}[style=codeblock,language=JavaScript,caption={Interface Chat AuraCompanion (frontend)}]
const AuraCompanion = () => {
  const [open, setOpen]       = useState(false);
  const [messages, setMessages]= useState([]);
  const [input, setInput]     = useState('');
  const [loading, setLoading] = useState(false);
  const { token }             = useAuth();

  const sendMessage = async () => {
    const userMsg = { role: 'user', content: input };
    setMessages(prev => [...prev, userMsg]);
    setInput('');
    setLoading(true);
    try {
      const { data } = await axios.post('/api/companion/chat',
        { messages: [...messages, userMsg] },
        { headers: { Authorization: `Bearer ${token}` } });
      setMessages(prev => [
        ...prev, { role: 'assistant', content: data.reply }
      ]);
    } finally {
      setLoading(false);
    }
  };

  return (
    <div className="fixed bottom-6 right-6 z-50">
      <AnimatePresence>
        {open && (
          <motion.div
            initial={{ scale: 0, opacity: 0 }}
            animate={{ scale: 1, opacity: 1 }}
            exit={{ scale: 0, opacity: 0 }}
            className="w-80 h-[450px] bg-white dark:bg-gray-900
                       rounded-2xl shadow-2xl flex flex-col overflow-hidden"
          >
            <div className="bg-indigo-600 p-3 text-white font-bold
                            flex items-center gap-2">
              <Sparkles size={18} /> Melo -- Votre IA Maison Intelligente
            </div>
            <div className="flex-1 overflow-y-auto p-3 space-y-3">
              {messages.map((m, i) => (
                <div key={i} className={`flex
                  ${m.role === 'user' ? 'justify-end' : 'justify-start'}`}>
                  <div className={`rounded-xl px-3 py-2 max-w-[80%] text-sm
                    ${m.role === 'user'
                      ? 'bg-indigo-100'
                      : 'bg-gray-100 dark:bg-gray-800'}`}>
                    {m.content
                      .replace(/\[ACTION:.*?\]/g, '')
                      .replace(/\[SCENE:.*?\]/g, '')}
                  </div>
                </div>
              ))}
              {loading && (
                <div className="text-xs text-gray-400 animate-pulse">
                  Melo reflechit...
                </div>
              )}
            </div>
            <div className="p-3 border-t flex gap-2">
              <input value={input} onChange={e => setInput(e.target.value)}
                onKeyDown={e => e.key === 'Enter' && sendMessage()}
                className="flex-1 rounded-lg border px-3 py-1.5 text-sm"
                placeholder="Posez une question a Melo..." />
              <button onClick={sendMessage}
                className="bg-indigo-600 text-white rounded-lg px-3 py-1.5 text-sm">
                Envoyer
              </button>
            </div>
          </motion.div>
        )}
      </AnimatePresence>
      <button onClick={() => setOpen(!open)}
        className="w-14 h-14 rounded-full bg-indigo-600 text-white shadow-lg
                   flex items-center justify-center hover:bg-indigo-700 transition">
        <Bot size={28} />
      </button>
    </div>
  );
};
\end{lstlisting}

\section{Modèle de Scénario et Moteur d'Automatisation}

\subsection{Schéma Scénario}

\begin{lstlisting}[style=codeblock,language=JavaScript,caption={Sch\'ema Sc\'enario (backend/models/Scenario.js)}]
const scenarioSchema = new mongoose.Schema({
  name:         { type: String, required: true },
  description:  { type: String },
  houseCode:    { type: String, required: true },
  isActive:     { type: Boolean, default: true },
  triggerType:  { type: String, enum: ['time'], default: 'time' },
  scheduleTime: { type: String, required: true },  // Format "HH:MM"
  action:       { type: String, required: true },  // ex. "mood:Morning"
  targetDevice: { type: mongoose.Schema.Types.ObjectId, ref: 'Device' },
  lastTriggeredAt: { type: Date }
}, { timestamps: true });
\end{lstlisting}

\subsection{Service Planificateur}

\begin{lstlisting}[style=codeblock,language=JavaScript,caption={Planificateur de Sc\'enarios (backend/scheduler.js)}]
const { schedule } = require('node-cron');

const initScheduler = (mqttClient, io) => {
  // Executer chaque minute
  schedule('* * * * *', async () => {
    const now = new Date();
    const currentTime = `${String(now.getHours()).padStart(2,'0')}:`
                      + `${String(now.getMinutes()).padStart(2,'0')}`;

    const scenarios = await Scenario.find({
      isActive: true,
      scheduleTime: currentTime
    });

    for (const scenario of scenarios) {
      try {
        if (scenario.action.startsWith('mood:')) {
          const scene = scenario.action.split(':')[1];
          await activateScene(scene, scenario.houseCode, mqttClient, io);
        } else if (scenario.action === 'lock_all') {
          await lockAllDoors(scenario.houseCode, mqttClient);
        } else if (scenario.action.startsWith('toggle:light:')) {
          const state = scenario.action.endsWith(':on') ? 'ON' : 'OFF';
          await toggleLights(scenario.houseCode, state, mqttClient);
        }
        scenario.lastTriggeredAt = new Date();
        await scenario.save();

        io.to(scenario.houseCode).emit('scenario:triggered', {
          name: scenario.name, action: scenario.action
        });

        await AuditLog.create({
          actorName: 'Planificateur',
          actorRole: 'system',
          action: `Scenario "${scenario.name}" declenche : ${scenario.action}`,
          category: 'scenario',
          houseCode: scenario.houseCode
        });
      } catch (err) {
        console.error(`[Planificateur] Erreur scenario ${scenario._id}:`, err);
      }
    }
  });
};
\end{lstlisting}

\section{Surveillance Énergétique}

\subsection{Algorithme de Calcul de la Consommation}

La consommation énergétique est estimée à partir de la collection AuditLog. Chaque
événement ON pour un appareil marque le début d'une période de consommation. L'événement
OFF correspondant (ou la fin de session) la ferme. Le kWh pour chaque période est :
\[ \text{kWh} = \frac{W_{\text{appareil}} \times \Delta t_{\text{heures}}}{1000} \]

\begin{lstlisting}[style=codeblock,language=JavaScript,caption={Endpoint \'Energie (backend/routes/energy.js)}]
router.get('/hourly', protect, async (req, res) => {
  const WATTAGE = { light:10, lamp:10, 'rgb-light':10, ac:1500,
                    tv:80, fan:50, outlet:200, 'washing-machine':500 };

  const logs = await AuditLog.find({
    houseCode: req.user.houseCode,
    category:  'device',
    createdAt: { $gte: new Date(Date.now() - 86400000) }  // Dernieres 24h
  }).sort({ createdAt: 1 });

  // Construire la carte de consommation horaire
  const hourly = Array(24).fill(0);
  const openPeriods = {};  // deviceId -> { startHour, watt }

  for (const log of logs) {
    const deviceId = log.metadata?.deviceId;
    const action   = log.action;
    const hour     = new Date(log.createdAt).getHours();
    const watt     = WATTAGE[log.metadata?.deviceType] ?? 10;

    if (action.includes('ON') || action.includes('OPEN')) {
      openPeriods[deviceId] = { startHour: hour, watt };
    } else if (openPeriods[deviceId]) {
      const duration = (hour - openPeriods[deviceId].startHour) || 1;
      hourly[hour] += (openPeriods[deviceId].watt * duration) / 1000;
      delete openPeriods[deviceId];
    }
  }

  res.json({ hourly, totalKwh: hourly.reduce((a,b) => a+b, 0) });
});
\end{lstlisting}

\section{Système de Demandes de Permission}

\subsection{Modèle DemandePermission}

\begin{lstlisting}[style=codeblock,language=JavaScript,caption={Sch\'ema DemandePermission (backend/models/PermissionRequest.js)}]
const permissionRequestSchema = new mongoose.Schema({
  requester:   { type: mongoose.Schema.Types.ObjectId, ref: 'User', required: true },
  houseAdmin:  { type: mongoose.Schema.Types.ObjectId, ref: 'User', required: true },
  actionKey:   { type: String, required: true },   // ex. 'room:kitchen'
  actionLabel: { type: String, required: true },   // libelle lisible
  room:        { type: String },                   // piece demandee (optionnel)
  status:      { type: String,
                 enum: ['pending','approved','denied'], default: 'pending' },
  reviewer:    { type: mongoose.Schema.Types.ObjectId, ref: 'User' },
  reviewNotes: { type: String },
  adminReadAt:     { type: Date },
  requesterReadAt: { type: Date }
}, { timestamps: true });
\end{lstlisting}

\subsection{Flux d'Approbation de Permission}

\begin{lstlisting}[style=codeblock,language=JavaScript,caption={Approbation Demande de Permission (backend/routes/permissions.js)}]
router.post('/:id/approve', protect, admin, async (req, res) => {
  const { reviewNotes } = req.body;
  const request = await PermissionRequest.findById(req.params.id)
    .populate('requester');

  if (!request || request.status !== 'pending')
    return res.status(400).json({ message: 'Demande introuvable ou deja traitee' });

  request.status      = 'approved';
  request.reviewer    = req.user._id;
  request.reviewNotes = reviewNotes;
  await request.save();

  // Accorder la permission au resident
  const resident = await User.findById(request.requester._id);
  if (!resident.permissions.includes(request.actionKey)) {
    resident.permissions.push(request.actionKey);
  }
  if (request.room && !resident.room) {
    resident.room = request.room;
  }
  await resident.save();

  // Notifier le resident en temps reel
  io.to(`user:${resident._id}`).emit('permission:updated', {
    status: 'approved', actionKey: request.actionKey
  });

  res.json({ success: true });
});
\end{lstlisting}

\section{Revue et Rétrospective Sprint 5}

\begin{tcolorbox}[goalbox,title={Revue Sprint 5}]
\textbf{Complété :} 4 stories sur 4 (34 PS). La démo Melo a impressionné. \\
\textbf{Points forts :} Taper ``Allume toutes les lumières de la chambre à 80\%'' a fait
illuminer la LED RGB chambre à 80\% de luminosité en 1,5 secondes. Le graphique d'énergie
a montré des pics de consommation réalistes lors du déclenchement du scénario climatisation. \\
\textbf{Rétrospective :}
\begin{itemize}[leftmargin=1cm,itemsep=2pt]
  \item \textbf{Ce qui a bien fonctionné :} La latence d'inférence Groq ($<$800~ms) rend
        Melo très réactif.
  \item \textbf{Défi :} L'analyse des blocs ACTION a nécessité des regex soigneuses pour
        éviter les faux positifs dans les réponses conversationnelles.
  \item \textbf{Amélioration :} Ajout de \texttt{temperature: 0.3} pour réduire les noms
        d'appareils hallucités dans les réponses LLM.
\end{itemize}
\end{tcolorbox}
